% ============================================================
% SGCV-IA — Manuscrito para Requirements Engineering (Springer)
% Plantilla oficial: sn-jnl.cls (Springer Nature LaTeX template)
% Enfoque 3 — Explicabilidad como Requisito No Funcional
% ============================================================

\documentclass[pdflatex,sn-mathphys-num,Numbered]{sn-jnl}

\usepackage{graphicx}
\usepackage{multirow}
\usepackage{amsmath,amssymb,amsfonts}
\usepackage{booktabs}
\usepackage{manyfoot}
\usepackage{xcolor}


\raggedbottom

\begin{document}

% ============================================================
% TITULO, AUTORIA, FILIACION
% ============================================================

\title[SGCV-IA: Explainability as a non-functional requirement]{Explainability as a non-functional requirement in an AI-assisted veterinary clinical system: a case study}

%%=============================================================
%% Autoria
%%=============================================================
\author*[1]{\fnm{Alberto Jeanpool} \sur{Marcillo Ponce}}\email{amarcillop@uteq.edu.ec}
\author[1]{\fnm{Robyn Willian} \sur{Amagua Sacón}}\email{ramaguas@uteq.edu.ec}
\author[1]{\fnm{Anthony Alfredo} \sur{Vera Gómez}}\email{averag10@uteq.edu.ec}
\author[1]{\fnm{Jhon Alexander} \sur{Mesías Quijije}}\email{jmesiasq@uteq.edu.ec}
\author[1]{\fnm{Carlos Daniel} \sur{Barrionuevo Fuentes}}\email{cbarrionuevof@uteq.edu.ec}
\author[1]{\fnm{Gleiston Ciceron} \sur{Guerrero Ulloa}}\email{gguerrero@uteq.edu.ec}

\affil*[1]{\orgdiv{Facultad de Ciencias de la Computación}, \orgname{Universidad Técnica Estatal de Quevedo (UTEQ)}, \orgaddress{\city{Quevedo}, \country{Ecuador}}}

% ============================================================
% RESUMEN ESTRUCTURADO — Contexto/Objetivo/Método/Resultados/Conclusiones
% PENDIENTE: completar Resultados y Conclusiones cuando existan datos
% ============================================================

\abstract{
\textbf{Context.} TODO — situar el problema: sistemas clínicos veterinarios con módulos de diagnóstico asistido por IA carecen frecuentemente de explicabilidad operacionalizada como requisito verificable. \\
\textbf{Objective.} TODO — determinar qué requisitos de explicabilidad, operacionalizados como RNF verificables, resultan necesarios y suficientes para los distintos perfiles de usuarios (técnicos y no técnicos) de un sistema clínico veterinario con IA. \\
\textbf{Method.} Applied single-case study grounded in Requirements Engineering. Requirements were elicited through semi-structured interviews and direct observation, analyzed qualitatively, operationalized as functional and non-functional requirements, and linked through requirements traceability. Explainability was treated as a non-functional requirement of the AI-assisted diagnostic functionality. \\
\textbf{Results.} [PENDIENTE — no completar hasta que existan datos analizados por los scripts de \texttt{06\_Experimento/scripts\_analisis/}]. \\
\textbf{Conclusions.} [PENDIENTE — depende de Results].
}

\keywords{Explainability, Non-functional requirements, AI-assisted diagnosis, Veterinary informatics, User studies}

\maketitle

% ============================================================
% 1. INTRODUCCION
% ============================================================

\section{Introduction}\label{sec1}

Artificial intelligence (AI)-assisted diagnostic support is increasingly embedded into clinical management systems, including veterinary practice, where automated recommendations directly influence treatment decisions taken by veterinary professionals and are also relevant to non-technical administrative staff who manage clinical workflows and reporting around those decisions. The SGCV-IA system --- a real-world veterinary clinic management system developed for a client organization in the province of Los R\'ios, Ecuador --- incorporates an AI-based diagnostic module whose outputs must be trusted and actionable by users with markedly different technical backgrounds. Studying Requirements Engineering in this context matters because, in AI-assisted clinical settings, the absence of explicit and verifiable explainability requirements can directly undermine user trust and, ultimately, patient (animal) safety.

Despite growing attention to explainability in software engineering, most existing work operationalizes explainability as a general software quality attribute rather than as a set of verifiable, stakeholder-specific non-functional requirements (NFRs) elicited through empirical fieldwork. Chazette et al.~\cite{chazette2021} propose a definition, model, and knowledge catalogue for explainability as a software quality concern, and Chazette et al.~\cite{chazette2022} extend this into an end-to-end framework spanning requirements analysis to system evaluation. However, empirical validation of this framework in real-world, resource-constrained, Latin American clinical domains --- and specifically in veterinary informatics --- remains scarce, leaving a gap in understanding how explainability requirements should be operationalized and validated for heterogeneous user profiles (technical vs.\ non-technical) in such settings.

This work contributes: (1) an empirical characterization of explainability concerns for a real-world AI-assisted veterinary diagnostic system; (2) their operationalization as explicit and verifiable non-functional requirements; and (3) traceability connecting stakeholder evidence, requirements, and system functionality.

Accordingly, this study addresses two research questions: \textbf{RQ1} asks what types of explanations users of the SGCV-IA system require to appropriately understand and assess the automated suggestions produced by the AI-assisted diagnostic component; \textbf{RQ2} asks to what extent these explainability concerns can be operationalized as explicit and verifiable non-functional requirements within the SGCV-IA requirements specification.

% ============================================================
% 2. TRABAJO RELACIONADO (Sección 5.3 — esquema Kitchenham y Charters)
% ============================================================

\section{Related Work}\label{sec2}

The related work was organized around three strands directly relevant to the present study: explainability as a software quality and requirements concern, Requirements Engineering for AI-based systems (RE4AI), and explainability in healthcare and other human-centred contexts. The selection was intentionally focused on studies that inform the operationalization, elicitation, specification, or evaluation of explainability and AI-related requirements. This is a focused literature review rather than a systematic literature review; consequently, no exhaustive database hit count is claimed. The comparison below contains 15 closely related studies selected for direct relevance to the research questions and methodological framing.

Research on explainability in software engineering has increasingly moved beyond model-level interpretability toward the consideration of explainability as a property that must be addressed during Requirements Engineering. Chazette and Schneider~\cite{chazette2020} characterize explainability as a non-functional requirement and discuss challenges associated with its elicitation and operationalization. Chazette et al.~\cite{chazette2021} provide a definition, model, and knowledge catalogue, while their subsequent work connects requirements analysis, design, and evaluation~\cite{chazette2022}. Deters et al.~\cite{deters2024} develop a quality model for explainability, and Droste et al.~\cite{droste2025} propose and evaluate an operational taxonomy of explanation needs. Together, these studies support treating explainability as an engineered and elicitable software concern rather than exclusively as a property of a machine-learning model.

A second strand concerns AI-based systems and the emergence of requirements challenges that differ from those of conventional software. Ahmad et al.~\cite{ahmad2023} identified 43 primary studies in a systematic mapping of RE4AI and reported limited empirical evaluation of ethics, explainability, and trust. Habiba et al.~\cite{habiba2024} subsequently mapped 126 primary studies and identified requirements specification, explainability, and the gap between machine-learning engineers and end users among the prevalent RE4AI challenges. Habibullah et al.~\cite{habibullah2023} show that machine-learning systems have quality demands that differ from traditional systems and examine how practitioners currently handle non-functional requirements. Heyn et al.~\cite{heyn2024} provide complementary empirical evidence from ten practitioner interviews concerning requirements for training data and runtime monitoring in critical ML systems. These studies indicate that AI-related requirements need to account for uncertainty, data, runtime behaviour, human stakeholders, and quality attributes beyond traditional functional behaviour.

A third strand focuses on users, healthcare, and the relationship between explanations and human judgement. Balasubramaniam et al.~\cite{balasubramaniam2023} connect transparency and explainability with requirements derived from ethical concerns, while de Brito Duarte et al.~\cite{debritoduarte2023} examine explainability in relation to warranted trust. Freyer et al.~\cite{freyer2024} synthesize ethical reasons for explainability in healthcare AI decision-support systems, and Mienye et al.~\cite{mienye2024} review explainable AI in healthcare and its adoption challenges. Rosenbacke et al.~\cite{rosenbacke2024} show that explanations can increase or decrease clinicians' trust depending on context. Habiba et al.~\cite{habiba2025} further report, from interviews with ML practitioners, diverse interpretations of explainability and recurring difficulties in communicating explainability needs to non-technical stakeholders.

The literature also provides methodological support for making requirements explicit and verifiable. Großer et al.~\cite{grosser2024} empirically compare requirement-template systems and show that structured requirement formulations can affect quality characteristics such as ambiguity and readability. Umar and Lano~\cite{umar2024} review automated support for Requirements Engineering and show that human intervention remains common across automated RE activities. Finkbeiner and Siber~\cite{finkbeiner2025} address explainability requirements from a formal perspective by modelling them as hyperproperties, illustrating that explainability can be formulated as a system property that is subject to verification.

Veterinary informatics provides a comparatively less developed application domain for this line of research. Lustgarten et al.~\cite{lustgarten2020} describe veterinary informatics as a field connecting veterinary practice with computational and biomedical methods. Although veterinary AI research has expanded, the literature identified here does not provide a comparable body of empirical Requirements Engineering studies focused specifically on explainability requirements in veterinary clinical-management systems.

\begin{table}[tp]
\caption{Comparison of closely related studies}
\label{tab1}
\centering
\tiny
\setlength{\tabcolsep}{1.5pt}
\renewcommand{\arraystretch}{0.85}
\begin{tabular}{@{}p{2.4cm}p{0.7cm}p{2.0cm}p{2.0cm}p{3.2cm}@{}}
\toprule
Reference & Year & Focus & Study type & Relation to the present study \\
\midrule
Chazette and Schneider~\cite{chazette2020} & 2020 & Explainability as NFR & Conceptual study & Establishes explainability as a non-functional requirement. \\
Chazette et al.~\cite{chazette2021} & 2021 & Explainability model & Framework study & Provides a conceptual model and knowledge catalogue for explainability. \\
Chazette et al.~\cite{chazette2022} & 2022 & Requirements-to-evaluation & Framework study & Connects requirements analysis with explainability evaluation. \\
Ahmad et al.~\cite{ahmad2023} & 2023 & RE for AI systems & Systematic mapping & Identifies gaps in empirical RE4AI research, including explainability and trust. \\
Habibullah et al.~\cite{habibullah2023} & 2023 & ML NFRs & Empirical study & Examines practitioner challenges in specifying NFRs for ML systems. \\
Balasubramaniam et al.~\cite{balasubramaniam2023} & 2023 & Transparency and explainability & Conceptual study & Links ethical transparency concerns with software requirements. \\
Freyer et al.~\cite{freyer2024} & 2024 & Healthcare explainability & Systematic review & Synthesizes ethical reasons for explainability in healthcare AI. \\
Mienye et al.~\cite{mienye2024} & 2024 & XAI in healthcare & Review & Surveys healthcare XAI and adoption challenges. \\
Habiba et al.~\cite{habiba2024} & 2024 & RE4AI maturity & Systematic mapping & Identifies explainability and stakeholder gaps in RE4AI. \\
Heyn et al.~\cite{heyn2024} & 2024 & ML requirements & Empirical interview study & Shows challenges in specifying data and runtime requirements for critical ML systems. \\
Deters et al.~\cite{deters2024} & 2024 & Explainability quality model & REFSQ study & Provides a quality-model perspective for eliciting explainability needs. \\
Rosenbacke et al.~\cite{rosenbacke2024} & 2024 & XAI and clinician trust & Systematic review & Shows that explanation effects on trust are context-dependent. \\
Droste et al.~\cite{droste2025} & 2025 & Explainability needs taxonomy & Empirical study & Provides an operational taxonomy to support explainability elicitation. \\
Habiba et al.~\cite{habiba2025} & 2025 & Practitioner explainability & Interview study & Highlights diverse interpretations and communication challenges with non-technical stakeholders. \\
Sporsem et al.~\cite{sporsem2026} & 2026 & Clinical explainability & Longitudinal case study & Provides a close clinical precedent for operationalizing explainability requirements. \\
\bottomrule
\end{tabular}
\end{table}

Taken together, the reviewed literature supports four observations. First, explainability is increasingly recognized as a software-engineering concern that should be addressed during Requirements Engineering rather than exclusively after model development~\cite{chazette2020,chazette2021,chazette2022,deters2024,droste2025}. Second, AI-based systems introduce requirements challenges involving explainability, trust, ethics, uncertainty, data, and the relationship between technical developers and end users~\cite{ahmad2023,habiba2024,habibullah2023,heyn2024}. Third, explainability is stakeholder- and context-dependent, and explanations do not uniformly produce appropriate trust~\cite{debritoduarte2023,freyer2024,rosenbacke2024,habiba2025}. Fourth, veterinary informatics remains comparatively underexplored for empirical Requirements Engineering studies focused on explainability~\cite{lustgarten2020}.

The resulting research gap is therefore not simply the absence of veterinary research on artificial intelligence. Existing literature already provides conceptual frameworks, taxonomies, quality models, and healthcare evidence concerning explainability. The unresolved problem is the empirical integration of these perspectives into a Requirements Engineering process in which stakeholder evidence from a real veterinary organization is transformed into explicit and verifiable explainability non-functional requirements and traced to the corresponding system functionality.

The present study addresses this gap through the SGCV-IA case. In particular, it examines explainability requirements associated with the AI-assisted diagnostic functionality and their relationship with RNF-18 and use case CU-06. The study therefore contributes an empirical Requirements Engineering perspective on explainability in a veterinary clinical-management context.

\section{Methodology}\label{sec3}

This study adopts an applied single-case design grounded in Requirements Engineering. The case is the SGCV-IA system, a veterinary clinical-management system incorporating an AI-assisted diagnostic functionality. The methodological focus of the present study is the elicitation, analysis, specification, and traceability of requirements, with particular attention to explainability as a non-functional requirement of the AI-assisted diagnostic component.

\subsection{Research questions and PICOC}\label{subsec1}

The study is guided by the following research questions:

\begin{itemize}
\item \textbf{RQ1.} What types of explanations do users of the SGCV-IA system require to appropriately understand and assess the automated suggestions produced by the AI-assisted diagnostic component?
\item \textbf{RQ2.} To what extent can the identified explainability concerns be operationalized as explicit and verifiable non-functional requirements within the SGCV-IA requirements specification?
\end{itemize}

The research questions are structured according to PICOC:

\begin{itemize}
\item \textbf{Population (P):} stakeholders involved in the clinical and organizational context of the veterinary system, including veterinary professionals and other relevant users identified during requirements elicitation.
\item \textbf{Intervention (I):} the elicitation and operationalization of explainability concerns as explicit non-functional requirements for the AI-assisted diagnostic functionality.
\item \textbf{Comparison (C):} comparison of stakeholder perspectives and requirements evidence where differences or convergences are identified in the empirical material.
\item \textbf{Outcome (O):} identified explainability concerns, their transformation into verifiable requirements, and their traceability to stakeholder evidence and system functionality.
\item \textbf{Context (C):} the real-world organizational context associated with a veterinary clinic in Ecuador and the development of the SGCV-IA system.
\end{itemize}

\subsection{Case study design}\label{subsec2}

The study follows the case study approach described by Runeson and H\"ost~\cite{runeson2009}. SGCV-IA constitutes the single case under investigation. The analysis is embedded in the Requirements Engineering process and considers stakeholder evidence, requirements artefacts, and the resulting traceability relationships.

The case was investigated in the context of real veterinary clinical and administrative activities. The empirical material was collected during the requirements-elicitation stage and subsequently incorporated into the requirements specification. The study does not treat the AI component as an isolated machine-learning artefact; instead, the AI-assisted diagnostic functionality is analyzed as part of a socio-technical software system in which professional users remain responsible for interpreting and acting upon the generated suggestions.

\subsection{Participants and recruitment}\label{subsec3}

The requirements-elicitation evidence identifies participants using anonymized identifiers P01--P16. The project repository maintains the correspondence between participant identifiers and real identities exclusively within a restricted area; the public project documentation does not reproduce this correspondence. This anonymization procedure is explicitly documented in the final ERS.

Participants were selected from the organizational context relevant to the veterinary system and contributed information concerning clinical workflows, administrative activities, information needs, and expectations regarding AI-assisted diagnostic support. The documented evidence includes interviews with veterinary professionals and observations of activities performed in the veterinary clinical setting.

Participation was voluntary and based on informed consent. Identifying information and original multimedia materials are handled separately from the anonymized research artefacts. The final ERS specifies that signed consent materials containing identifying information are stored in the restricted repository area, while anonymized or masked copies are used for dissemination.

\subsection{Data collection}\label{subsec4}

The principal data-collection techniques documented in the final ERS were semi-structured interviews and direct observation.

\textbf{Semi-structured interviews.} Interviews were conducted with veterinary professionals to identify current workflows, operational problems, information needs, expectations regarding the system, and conditions associated with the use of AI-assisted diagnostic suggestions. The interview instrument and corresponding evidence are included in the project repository. The documented evidence includes interview records associated with anonymized participant identifiers.

\textbf{Direct observation.} Observation was used to complement interview evidence by examining clinical and administrative activities in their organizational context. The final ERS records observation evidence concerning patient registration, clinical information management, and other processes relevant to the system requirements.

\textbf{Documentary evidence.} Organizational documents and other project artefacts were reviewed when necessary to understand processes and support requirements specification. These materials complement, rather than replace, the primary field evidence.

No questionnaire or survey is treated as a requirements-elicitation instrument in this stage of the study. The final ERS explicitly states that requirements information for this stage was obtained through semi-structured interviews and observation.

\subsection{Requirements analysis}\label{subsec5}

The collected evidence was analyzed to identify stakeholder needs and transform them into structured requirements. The process consisted of four main activities: (1) identification of relevant statements and needs in the empirical evidence; (2) abstraction and consolidation of recurring needs; (3) formulation of functional and non-functional requirements; and (4) traceability of the resulting requirements to their evidence and system artefacts.

Particular attention was given to explainability concerns associated with the AI-assisted diagnostic functionality. Explainability was treated as a software quality concern and subsequently operationalized as a non-functional requirement, following the Requirements Engineering perspective established by Chazette et al.~\cite{chazette2020,chazette2021,chazette2022}.

In the final requirements specification, RNF-18 addresses explainability of the AI-assisted diagnostic functionality. The requirement is associated with the diagnostic use case CU-06 and specifies information that must be presented to the veterinary professional together with the AI-generated suggestion. This traceability connects the explainability concern to an identifiable requirement and system function.

\subsection{Requirements traceability}\label{subsec6}

Traceability was used to maintain the relationship between empirical evidence, raw requirements, formal requirements, system functionality, and design artefacts. The final ERS reports an extended traceability matrix containing 60 rows and links requirements with their corresponding sources, objectives, normative considerations, design elements, and associated artefacts.

This approach reduces the risk that explainability requirements become detached from the stakeholder evidence from which they originated. It also provides a basis for subsequent verification because each formal requirement can be related to an identifiable source and to the functionality in which it is intended to be realized.

\subsection{Ethical and data-protection considerations}\label{subsec7}

The study involves human participants and therefore incorporates informed consent and confidentiality measures. Participant identities are represented by anonymized identifiers P01--P16 in the public project materials. The mapping between these identifiers and real identities is retained exclusively in a restricted repository location.

Original consent documents, audio, video, and other potentially identifying materials are separated from the public research artefacts. The project documentation states that the handling of personal data follows the requirements established under Ecuador's Ley Orgánica de Protección de Datos Personales~\cite{lopdp2021}.

No personally identifying participant information is reproduced in this manuscript.

\subsection{Data analysis and reproducibility}\label{subsec8}

The qualitative analysis is supported by interview transcripts, observation evidence, documentary material, and thematic coding. The project repository contains anonymized transcripts, coding materials, requirements artefacts, and traceability information. The final ERS documents the organization of these materials and distinguishes restricted evidence from materials intended for dissemination.

The terminal empirical validation and its statistical analysis are not reported in this manuscript version. Results, statistical tests, effect sizes, confidence intervals, and conclusions dependent on the experimental dataset will be incorporated only after the corresponding experimental evidence and analysis scripts have been completed and verified.

\subsection{Methodological limitations}\label{subsec9}

The single-case design provides detailed contextual evidence but does not by itself establish broad generalizability across veterinary organizations or AI-assisted clinical systems. The study therefore emphasizes contextualized requirements evidence and traceability rather than population-level statistical inference. The implications of this design for internal, external, construct, and conclusion validity are discussed in Section~\ref{sec6}.

% ============================================================
% 4. RESULTADOS — NO ESCRIBIR TODAVIA
% Advertencia expresa de la guía: la sección de Resultados NO puede
% redactarse antes de que existan los datos analizados por los scripts.
% ============================================================

\section{Results}\label{sec4}
% PENDING: to be completed only after the terminal experimental dataset
% and corresponding analysis scripts have been completed and verified.

% ============================================================
% 5. DISCUSION — NO ESCRIBIR TODAVIA (depende de Resultados)
% ============================================================

\section{Discussion}\label{sec5}
% NO COMPLETAR. Tres subapartados: implicaciones para la investigación,
% implicaciones para la práctica, hallazgos sorprendentes/contraintuitivos.

% ============================================================
% 6. AMENAZAS A LA VALIDEZ (C10 — apoyo de Barrionuevo)
% ============================================================

\section{Threats to Validity}\label{sec6}

The validity of this study is discussed according to the four categories commonly used in empirical software engineering: internal, external, construct, and conclusion validity~\cite{wohlin2012}. The discussion is adapted to the characteristics of a qualitative, single-case Requirements Engineering study and complements the case-study guidance of Runeson and H\"ost~\cite{runeson2009}.

\subsection{Internal validity}\label{tv:internal}

A first threat concerns researcher interpretation during the transformation of interview and observation evidence into formal requirements. Because the study involves qualitative analysis, the identification, abstraction, and consolidation of stakeholder needs may be influenced by the analysts' judgement. This threat was mitigated by maintaining traceability between empirical evidence, raw requirements, formal requirements, use cases, and design artefacts. The extended traceability matrix provides an auditable chain linking requirements to their documented sources. Nevertheless, interpretive decisions cannot be eliminated completely, and alternative analysts could classify or formulate some stakeholder concerns differently.

A second threat arises from the possibility that interview statements alone may not fully represent actual clinical and administrative practice. To reduce dependence on a single source of evidence, semi-structured interviews were complemented with direct observation and documentary evidence from the organizational context. This triangulation strengthens the consistency of the elicited requirements. However, the observations were conducted within the same organizational context as the interviews, so they do not provide an independent replication of the findings in another veterinary organization.

\subsection{External validity}\label{tv:external}

The main threat to external validity is the single-case design. SGCV-IA was investigated within a specific veterinary organizational context in Ecuador, and the stakeholders, workflows, constraints, and regulatory environment may differ from those of other veterinary clinics, countries, or AI-assisted clinical systems. The study therefore does not claim statistical generalization. Instead, the case, stakeholder roles, elicitation procedures, requirements artefacts, and traceability relationships are documented in detail to support analytical comparison with similar contexts, as recommended for software-engineering case studies~\cite{runeson2009}. The residual limitation is that transferability to other organizations must be assessed by future studies in additional settings.

A second threat concerns domain specificity. Explainability needs identified for an AI-assisted veterinary diagnostic functionality may not transfer directly to human healthcare, other high-stakes domains, or AI systems used for different tasks. The study mitigates this threat by grounding its contribution in Requirements Engineering concepts and by explicitly identifying the SGCV-IA context and the diagnostic use case to which RNF-18 is traced. Even so, empirical replication is required before broader claims can be made about explainability requirements across domains.

\subsection{Construct validity}\label{tv:construct}

A primary construct-validity threat is the operationalization of explainability. Explainability is a broad and context-dependent quality concern, and reducing it to a finite set of software requirements may omit dimensions relevant to particular stakeholders. To mitigate this threat, explainability was treated explicitly as a non-functional requirement following established Requirements Engineering work on explainability~\cite{chazette2020,chazette2021,chazette2022}. In SGCV-IA, the concern is operationalized through RNF-18 and traced to the AI-assisted diagnostic use case CU-06, providing a concrete connection between the abstract construct and system behaviour. The residual limitation is that one operationalization cannot be assumed to cover every possible interpretation of explainability.

A second threat is the potential conflation of explainability with related constructs such as transparency, trust, usability, or diagnostic accuracy. The present study reduces this risk by focusing the analysis on the information that must accompany an AI-generated suggestion and on its specification as a verifiable software requirement, rather than treating user trust or model performance as equivalent measures of explainability. Nevertheless, stakeholder statements may naturally combine these concerns, and qualitative coding cannot completely isolate them from one another.

\subsection{Conclusion validity}\label{tv:conclusion}

The principal threat to conclusion validity at the present stage is the qualitative and contextual nature of the available evidence. The conclusions that can be supported by this part of the study concern the identification, specification, and traceability of explainability requirements; they do not establish population-level effects or statistically significant differences between user groups. This threat is mitigated by restricting claims to the documented evidence and by preserving traceability from stakeholder material to the resulting requirements. Accordingly, no statistical inference is reported in this manuscript version.

A second threat concerns the incompleteness of the terminal empirical validation at the time of this manuscript version. Statistical tests, effect sizes, confidence intervals, and conclusions dependent on the experimental dataset are intentionally excluded until the corresponding data and analysis scripts have been completed and verified. This separation prevents unsupported quantitative conclusions, but it also means that the present study cannot yet make claims about the measured effectiveness, satisfaction, or comparative performance of the operationalized explainability requirements. Those claims remain outside the scope of the current manuscript version.

% ============================================================
% 7. CONCLUSIONES Y TRABAJO FUTURO — depende de Resultados
% ============================================================

\section{Conclusions and Future Work}\label{sec7}
% TODO. Un párrafo por RQ; un párrafo de contribuciones consolidadas;
% al menos 3 líneas concretas de trabajo futuro.

% ============================================================
% 8. DISPONIBILIDAD DE DATOS, USO DE IA Y AGRADECIMIENTOS
% ============================================================

\backmatter

\section*{Data and materials availability}
The materials supporting the requirements-elicitation and analysis reported in this manuscript are maintained in the SGCV-IA project repository. The available research artefacts include anonymized interview transcripts, thematic-coding materials, requirements artefacts, and traceability information. Original consent forms and multimedia materials containing potentially identifying information are retained in a restricted repository area and are not publicly released. Anonymized or masked copies are used for dissemination where appropriate. No public DOI or persistent identifier is currently assigned to the project repository; therefore, no repository identifier is reported here. Access to restricted materials is subject to the applicable ethical and data-protection conditions.

\section*{Use of AI-assisted technologies}
During preparation of this manuscript, an AI-assisted language model (ChatGPT, GPT-5.6 Luna) was used to assist with drafting, restructuring, and language refinement of manuscript text. The tool was not used to generate, alter, or interpret empirical data, statistical results, tables, or conclusions. The authors remain responsible for verifying the accuracy, provenance, and scientific integrity of all content and for the final version of the manuscript. No personal, confidential, or identifying participant information was intentionally included in prompts.

\section*{Acknowledgements}
The authors thank the veterinary professionals and other participants who contributed to the requirements-elicitation activities. Participant identities are not disclosed in this manuscript in accordance with the project's anonymization and data-protection procedures.

% ============================================================
% REFERENCIAS
% ============================================================

\bibliography{referencias}

% ============================================================
\end{document}
